\section{Design Principles}
\label{sec:principles}

\tool{} is opinionated. Each principle below corresponds to a failure that actually happened during
development; Section~\ref{sec:controls} lists them with their numbers. The principles are enforced in
code (as gates, refusals and tests) rather than left as recommendations in documentation.

\begin{description}[leftmargin=1.2em,font=\normalfont\bfseries]
\item[P1. Ground truth by construction.] Every probe is first validated on a case whose answer is
  known. The benchmark records every generative component of every synthetic series (trend order and
  scale, seasonal periods, AR coefficients, changepoints, anomalies, intermittency), so ``does the model
  represent seasonality?'' can be checked against the truth before it is asked of real data. Every
  controlled perturbation (detrending, deseasonalizing, level shifts and six others) is applied to
  series whose components are known, so its effect on the input is known as well.
\item[P2. Alignment is measured, not declared.] A token's time span is measured from the model's own
  impulse response and checked against what the adapter declares. Cross-model analysis happens only
  on a common window axis (Section~\ref{sec:alignment}). A model whose tokens are not time-localized is
  routed to behavioral analysis only, with the verdict rendered in the report.
\item[P3. A causal claim needs a null, a reach control and a same-seed baseline.] Every intervention
  is compared with a random intervention of matched size on the \emph{same} series. Before any
  intervention is trusted, patching a layer into itself must change the forecast by exactly $0.0$,
  and patching in a different layer's activations must change it by a nonzero amount. Sampled models
  are seeded identically in every pair of compared calls.
\item[P4. The series is the unit of inference.] Windows within a series are not independent. All
  confidence intervals and $p$-values resample series (paired, cluster or validation-series
  bootstraps), all $p$-values are floored at $1/n_{\text{boot}}$, and every correction family's
  attainable minimum is checked against its threshold.
\item[P5. Explore, then confirm once.] Everything computed on the dev split is exploratory.
  Hypotheses are frozen with the content hash of the artifact they came from, and the sealed private
  split is opened once. A second look is recorded as a repeated look, and a fresh epoch must be minted
  to confirm again.
\item[P6. Every claim is typed, and failures are loud.] Each finding carries an evidence class
  (\ev{behavioral}, \ev{descriptive}, \ev{geometric}, \ev{linear-translatable},
  \ev{causal within-model}, \ev{illustrative}, \ev{held-out confirmed}) and a caveat generated from
  that class. An unsupported capability is skipped with a reason rendered in the report. An
  unresolvable slice disables the analysis rather than guessing, and a broken assumption raises.
\item[P7. Accessible by default.] The deliverable is one self-contained HTML file that a person with no
  interpretability background can read. Every figure has a caption and a ``What does this mean?''
  explainer that states the figure's purpose, how to read it and its limitations. The headline
  scorecard is computed from artifacts by a pure function; no verdict in it is written by hand.
\end{description}

Two further engineering rules protect reproducibility. New behavior is opt-in, so that sealed corpora
regenerate bit-exactly (golden hashes are pinned in the benchmark tests). And every stage's
configuration is fingerprinted, so a cached artifact computed under a different configuration is
refused rather than silently reused.
